\documentclass[10pt,a4paper]{amsart}
\usepackage[margin=19mm]{geometry}
\usepackage{amsmath,amssymb,mathtools}
\usepackage[scr=rsfs,cal=euler]{mathalfa}
\usepackage{xfrac}
\usepackage[unicode,hidelinks,bookmarks=false]{hyperref}
\newcommand{\ie}{i.e.\ }
\newcommand{\cf}{cf.\ }
\newcommand{\resp}{respectively\ }
\newcommand{\Subsection}{\subsection}
\providecommand{\nobreakdash}{\nobreak\hbox{-}}
\setlength{\emergencystretch}{2em}
\multlinegap0pt
\usepackage{amssymb}
\usepackage{enumerate}
\usepackage{mathtools}
\usepackage[shortlabels]{enumitem}
\usepackage{tikz}
\newtheorem{prop}{Proposition}
\newtheorem{thm}{Theorem}
\title{Ordinarization Numbers of Numerical Semigroups}
\author{Sogol Cyrusian, Nathan Kaplan}
\date{Source: arXiv:2506.10222v2 (CC BY 4.0)}
\begin{document}
\maketitle

\noindent\emph{Source and licence.} Ordinarization Numbers of Numerical Semigroups. Version arXiv:2506.10222v2 (CC BY 4.0), distributed by arXiv under the Creative Commons Attribution 4.0 International licence, http://creativecommons.org/licenses/by/4.0/, as recorded in the arXiv licence metadata for this exact version and shown on the abstract page https://arxiv.org/abs/2506.10222v2. Full licence notice: \texttt{licenses/arxiv-2506.10222-CC-BY-4.0.txt}.\par\smallskip

\noindent\textit{Editorial scope.} Counted unit: the article's prop 2genO (source0.tex lines 608--613). The article's complete original proof of it is delivered with it (source0.tex lines 694--758, one \texttt{proof} environment of 65 lines, which the article places away from the statement). No mathematics is written, repaired or re-derived: the statement and the proof are the article's own lines, in the article's own notation, and no retained line is substituted or re-flowed.  Retained uncounted as the source-local context the proof needs: lines 621--634.  Nothing was omitted from any retained span, so the package carries no omission record.  No retained line contains a citation, so the wrapper carries no bibliography and no bibliography entry.  No retained line contains a figure environment, an image-inclusion command or drawing code, so no image is delivered and the package declares no asset dependency.  Editorial formatting only, in addition: an AMS \texttt{amsart} preamble that restates the article's own declarations and macros used by the retained lines, namely the environments \texttt{prop}, \texttt{thm} on separate counters, together with the standard packages those lines need.  The retained environments print in this document as Proposition 1, Theorem 1 in order of appearance, whereas the article prints the same items with the numbering of its own full document.  The complete native source of the article as distributed in the arXiv e-print for this exact version is bundled unchanged as \texttt{sources/179718/source0.tex}.
\begin{thm}\label{thm:ord_e2}
Let $S = \langle a,b \rangle$ where $2\le a < b$ and $\gcd(a,b) = 1$.
\begin{enumerate}
\item Suppose $a$ is odd.  We have
\[
0 \le r(S) - \left( \frac{ab-a-4}{8} - \frac{b+3}{8a}\right) \le \frac{a}{4}-\frac{1}{4a}.
\]

\item Suppose $a$ is even.  We have
\[
0 \le r(S) - \left(\frac{ab-a-4}{8} \right)\le \frac{a}{4}.
\]
\end{enumerate}
\end{thm}

\begin{prop}\label{2genO}
Let $S = \langle a,b\rangle$ where $2\le a < b$ and $\gcd(a,b) = 1$. Write $g = \frac{ab-a-b+1}{2}$.  Then
\[
r(S) =\left\lfloor\frac{g}{b}\right\rfloor +\sum\limits_{n=0}^{\lfloor\frac{g}{b}\rfloor}\left\lfloor \frac{g-nb}{a}\right\rfloor.
\]
\end{prop}

\begin{proof}[Proof of Theorem \ref{thm:ord_e2}]
To simplify the notation, we write $g = \frac{ab-a-b+1}{2} = \frac{(a-1)(b-1)}{2}$.

We note that 
\[
\frac{g}{b} = \frac{(a-1)(b-1)}{2b} = \frac{a-1}{2}-\frac{a-1}{2b}.
\]
Since $1 \le a-1 < b$, we see that 
\[
0 < \frac{a-1}{2b} < \frac{1}{2}.
\]
This implies 
\begin{equation}\label{eqn:floor}
\left\lfloor \frac{g}{b} \right\rfloor = \begin{cases} \frac{a-3}{2} & \text{ if } a \text{ is odd} \\
\frac{a-2}{2} & \text{ if } a \text{ is even} 
\end{cases}.
\end{equation}

In order to compute $r(S)$, we need to compute the sum in Proposition \ref{2genO}.  We have
\begin{equation}\label{eqn:sums}
\sum_{n=0}^{\lfloor \frac{g}{b}\rfloor} \left\lfloor \frac{g-nb}{a} \right\rfloor = \sum_{n=0}^{\lfloor \frac{g}{b}\rfloor} \left(\frac{g-nb}{a} \right) - \sum_{n=0}^{\lfloor\frac{g}{b}\rfloor} \left\{\frac{g-nb}{a} \right\}.
\end{equation}
We consider each of these two sums separately.  Because of the difference in the expression for $\lfloor \frac{g}{b} \rfloor$ it is helpful to divide our analysis into cases based on whether $a$ is even or odd.

Suppose that $a$ is odd.  We have
\begin{align}
\sum_{n=0}^{\lfloor \frac{g}{b}\rfloor} \left(\frac{g-nb}{a} \right) & =  \sum_{n=0}^{ \frac{a-3}{2}} \left(\frac{g-nb}{a} \right) = \sum_{n=0}^{ \frac{a-3}{2}} \frac{g}{a} - \sum_{n=0}^{ \frac{a-3}{2}} \frac{nb}{a}  =  \frac{g}{a}\left( \frac{a-1}{2}\right) - \frac{b}{a} \sum_{n=0}^{ \frac{a-3}{2}} n \nonumber\\
&  =   \frac{(a-1)(b-1)}{2a} \left( \frac{a-1}{2} \right)- \frac{b}{a} \frac{\left(\frac{a-3}{2}\right) \left(\frac{a-1}{2}\right)}{2} \nonumber \\
& =  \frac{(ab - 2a + b + 2)(a - 1)}{8a}. \label{eqn:main_term_odd}
\end{align}
A similar analysis when $a$ is even shows that 
\begin{equation}\label{eqn:main_term_even}
\sum_{n=0}^{\lfloor \frac{g}{b}\rfloor} \left(\frac{g-nb}{a} \right) = \frac{ab - 2a + 2}{8}.
\end{equation}

We now turn to the next expression in the formula for $r(S)$,
\[
\sum_{n=0}^{\lfloor\frac{g}{b}\rfloor} \left\{\frac{g-nb}{a} \right\}.
\]

We first suppose that $a$ is odd. For any integer $n$, we see that $\left\{\frac{g-nb}{a}\right\}$ is a rational number equal to one of $\frac{0}{a},\frac{1}{a},\ldots, \frac{a-1}{a}$. The particular value of $\left\{\frac{g-nb}{a}\right\}$ is determined by $g-nb \pmod{a}$.

Consider the $a$ integers given by $g-nb$ where $0 \le n \le a-1$.  These values of $nb$ are distinct modulo $a$. Therefore, the integers $g-nb$ are distinct modulo $a$ as well.  We see that the $\frac{a-1}{2}$ values of $\left\{\frac{g-nb}{a}\right\}$ when $0 \le n \le \frac{a-3}{2}$ give $\frac{a-1}{2}$ distinct values in the set $\left\{\frac{0}{a},\frac{1}{a},\ldots, \frac{a-1}{a}\right\}$.  The smallest that the sum of these terms can be is 
\[
\frac{0}{a} + \frac{1}{a} + \cdots + \frac{\left(\frac{a-3}{2}\right)}{a} = \frac{(a-3)(a-1)}{8a}.
\]
The largest that the sum of these terms can be is 
\[
\frac{a-1}{a} + \frac{a-2}{a} + \cdots + \frac{a-\left(\frac{a-1}{2}\right)}{a} = 
\frac{\frac{a+1}{2}}{a}+ \cdots + \frac{a-1}{a}.
\]
This is the same as the previous sum except that $\frac{a+1}{2a} = \frac{1}{2} + \frac{1}{2a}$ has been added to each term.  Therefore,  the largest this sum can be is 
\[
\frac{(a-3)(a-1)}{8a} + \left(\frac{a+1}{2a}\right) \left( \frac{a-1}{2}\right) = \frac{(3a-1)(a-1)}{8a}.
\]
Combining these observations shows that 
\begin{equation}\label{eqn:fractional_odd}
 \frac{(a-3)(a-1)}{8a} \le \sum_{n=0}^{\lfloor\frac{g}{b}\rfloor} \left\{\frac{g-nb}{a} \right\} \le  \frac{(a-3)(a-1)}{8a} + \frac{a^2-1}{4a}.
\end{equation}
A similar analysis when $a$ is even shows that
\begin{equation}\label{eqn:fractional_even}
\frac{a-2}{8} \le \sum_{n=0}^{\lfloor\frac{g}{b}\rfloor} \left\{\frac{g-nb}{a} \right\} \le \frac{a-2}{8} + \frac{a}{4}.
\end{equation}
Combining the results of equations \eqref{eqn:floor}, \eqref{eqn:sums}, \eqref{eqn:main_term_odd}, \eqref{eqn:main_term_even}, \eqref{eqn:fractional_odd}, and \eqref{eqn:fractional_even} completes the proof.
\end{proof}

\end{document}
